% GENERATED FILE. Edit canonical lessons, then run npm run generate:books.
% canonical-source-hash: fnv1a64:471d9473a9b69d17
% canonical-lessons: TA-C374-tertal, TA-C374-vetpalar, TA-C374-vakkali, TA-C374-tokuti, TA-C374-piraccaram

\chapter{Elections}
\label{ch:ta-elections}
\hlchaptermodality{374}

\begin{chapteropening}
\textbf{By the end of this chapter:} \emph{I can talk about elections, candidates and voting.}

Complete the last lesson of chapter 374: I can talk about elections, candidates and voting.
\end{chapteropening}

\section[\texorpdfstring{\ta{தேர்தல்} (tērtal)}{தேர்தல் (tērtal)}]{\ta{தேர்தல்} (tērtal) — an election}

\label{lesson:TA-C374-tertal}



\begin{quote}
Before the new one: say the Tamil for a governor, then the Tamil for a parliament.
\end{quote}

\subsection*{You'll want to know: \ta{தேர்தல்}}
\textbf{\ta{தேர்தல்}} — \emph{tērtal} — \textquotedblleft{}an election\textquotedblright{}.

\begin{quote}
\ta{அடுத்த} \ta{மாதம்} \ta{தேர்தல்} \ta{நடக்கும்}.
\end{quote}

\begin{quote}
\emph{aṭutta mātam tērtal naṭakkum.}
\end{quote}

\begin{quote}
The election will be held next month.
\end{quote}

From \ta{தேர்}, to choose, as in \ta{தேர்ந்தெடு}, + -\ta{தல்}, the verbal noun.

\begin{cousinweb}
\textbf{\ta{தேர்தல்}} (\emph{tērtal}) is native: the verbal noun in -\ta{தல்} of \emph{tēr}, \textquotedblleft{}to examine, consider, choose\textquotedblright{} (DEDR 3419), the verb inside \ta{தேர்ந்தெடு}.
\end{cousinweb}

\begin{grammarlens}[title={What you've built}]
One more word for giving reasons and opinions.
\end{grammarlens}

\subsection*{Your turn}
\begin{itemize}
\raggedright
  \item \emph{Say it:} \emph{tērtal}
  \item \emph{Say it:} \emph{tērtal}, once more
  \item \emph{Say it:} \emph{aṭutta mātam tērtal naṭakkum.} — The election will be held next month.
  \item \emph{Recall:} say the Tamil for a governor, then the Tamil for a parliament, then say \emph{tērtal} again
\end{itemize}

\begin{tcolorbox}[breakable,colback=teal!4,colframe=teal!35!black,title={Before you move on}]
What does \ta{தேர்தல்} mean? (\textquotedblleft{}An election\textquotedblright{}.) Say it once more.
\end{tcolorbox}
\section[\texorpdfstring{\ta{வேட்பாளர்} (vēṭpāḷar)}{வேட்பாளர் (vēṭpāḷar)}]{\ta{வேட்பாளர்} (vēṭpāḷar) — a candidate}

\label{lesson:TA-C374-vetpalar}



\begin{quote}
Before the new one: say the Tamil for a parliament, then the Tamil for an election.
\end{quote}

\subsection*{You'll want to know: \ta{வேட்பாளர்}}
\textbf{\ta{வேட்பாளர்}} — \emph{vēṭpāḷar} — \textquotedblleft{}a candidate\textquotedblright{}.

\begin{quote}
\ta{இந்தத்} \ta{தேர்தலில்} \ta{ஐந்து} \ta{வேட்பாளர்கள்} \ta{போட்டியிடுகிறார்கள்}.
\end{quote}

\begin{quote}
\emph{intat tērtalil aintu vēṭpāḷarkaḷ pōṭṭiyiṭukiṟārkaḷ.}
\end{quote}

\begin{quote}
Five candidates are standing in this election.
\end{quote}

\ta{வேட்பு}, desire + -\ta{ஆளர்}, the one who: one who seeks the seat.

\begin{cousinweb}
\textbf{\ta{வேட்பாளர்}} (\emph{vēṭpāḷar}) is native: \emph{vēṭpu}, \textquotedblleft{}desire\textquotedblright{}, from \emph{vēḷ}, \textquotedblleft{}to desire\textquotedblright{} (DEDR 5528), with \emph{-āḷar}, \textquotedblleft{}one who\textquotedblright{}: one who seeks.
\end{cousinweb}

\begin{grammarlens}[title={What you've built}]
One more word for giving reasons and opinions.
\end{grammarlens}

\subsection*{Your turn}
\begin{itemize}
\raggedright
  \item \emph{Say it:} \emph{vēṭpāḷar}
  \item \emph{Say it:} \emph{vēṭpāḷar}, once more
  \item \emph{Say it:} \emph{intat tērtalil aintu vēṭpāḷarkaḷ pōṭṭiyiṭukiṟārkaḷ.} — Five candidates are standing in this election.
  \item \emph{Recall:} say the Tamil for a parliament, then the Tamil for an election, then say \emph{vēṭpāḷar} again
\end{itemize}

\begin{tcolorbox}[breakable,colback=teal!4,colframe=teal!35!black,title={Before you move on}]
What does \ta{வேட்பாளர்} mean? (\textquotedblleft{}A candidate\textquotedblright{}.) Say it once more.
\end{tcolorbox}
\section[\texorpdfstring{\ta{வாக்களி} (vākkaḷi)}{வாக்களி (vākkaḷi)}]{\ta{வாக்களி} (vākkaḷi) — to vote}

\label{lesson:TA-C374-vakkali}



\begin{quote}
Before the new one: say the Tamil for an election, then the Tamil for a candidate.
\end{quote}

\subsection*{You'll want to know: \ta{வாக்களி}}
\textbf{\ta{வாக்களி}} — \emph{vākkaḷi} — \textquotedblleft{}to vote\textquotedblright{}.

\begin{quote}
\ta{நான்} \ta{முதல்} \ta{முறையாக} \ta{வாக்களித்தேன்}.
\end{quote}

\begin{quote}
\emph{nāṉ mutal muṟaiyāka vākkaḷittēṉ.}
\end{quote}

\begin{quote}
I voted for the first time.
\end{quote}

Past (I …): \ta{வாக்களித்தேன்}, \emph{vākkaḷittēṉ}.

For whom: X-\ta{க்கு} \ta{வாக்களி}. \ta{வாக்கு}, a vote; \ta{வாக்காளர்}, a voter.

\begin{cousinweb}
\textbf{\ta{வாக்களி}} (\emph{vākkaḷi}) is built from \emph{vākku}, \textquotedblleft{}word, vote\textquotedblright{}, borrowed from Sanskrit \emph{vāc}, \textquotedblleft{}speech, word\textquotedblright{} (Monier-Williams), and native \emph{aḷi}, \textquotedblleft{}to show favour, to grant\textquotedblright{} (DEDR 301, where Malayalam \emph{aḷikkuka} means \textquotedblleft{}to bestow\textquotedblright{}): to grant one's word.
\end{cousinweb}

\begin{grammarlens}[title={What you've built}]
One more word for giving reasons and opinions.
\end{grammarlens}

\subsection*{Your turn}
\begin{itemize}
\raggedright
  \item \emph{Say it:} \emph{vākkaḷi}
  \item \emph{Say it:} \emph{vākkaḷi}, once more
  \item \emph{Say it:} \emph{nāṉ mutal muṟaiyāka vākkaḷittēṉ.} — I voted for the first time.
  \item \emph{Recall:} say the Tamil for an election, then the Tamil for a candidate, then say \emph{vākkaḷi} again
\end{itemize}

\begin{tcolorbox}[breakable,colback=teal!4,colframe=teal!35!black,title={Before you move on}]
What does \ta{வாக்களி} mean? (\textquotedblleft{}To vote\textquotedblright{}.) Say it once more.
\end{tcolorbox}
\section[\texorpdfstring{\ta{தொகுதி} (tokuti)}{தொகுதி (tokuti)}]{\ta{தொகுதி} (tokuti) — a constituency; a volume}

\label{lesson:TA-C374-tokuti}



\begin{quote}
Before the new one: say the Tamil for a candidate, then the Tamil for to vote.
\end{quote}

\subsection*{You'll want to know: \ta{தொகுதி}}
\textbf{\ta{தொகுதி}} — \emph{tokuti} — \textquotedblleft{}a constituency; a volume\textquotedblright{}.

\begin{quote}
\ta{எங்கள்} \ta{தொகுதிக்கு} \ta{ஒரு} \ta{புதிய} \ta{மருத்துவமனை} \ta{வேண்டும்}.
\end{quote}

\begin{quote}
\emph{eṅkaḷ tokutikku oru putiya maruttuvamaṉai vēṇṭum.}
\end{quote}

\begin{quote}
Our constituency needs a new hospital.
\end{quote}

Also a volume of a book: \ta{முதல்} \ta{தொகுதி}, the first volume.

\begin{cousinweb}
\textbf{\ta{தொகுதி}} (\emph{tokuti}) is native Dravidian: DEDR 3476 lists Tamil \emph{toku}, \textquotedblleft{}to assemble, collect\textquotedblright{}, with \emph{tokuti}, \textquotedblleft{}collection, group\textquotedblright{}.
\end{cousinweb}

\begin{grammarlens}[title={What you've built}]
One more word for giving reasons and opinions.
\end{grammarlens}

\subsection*{Your turn}
\begin{itemize}
\raggedright
  \item \emph{Say it:} \emph{tokuti}
  \item \emph{Say it:} \emph{tokuti}, once more
  \item \emph{Say it:} \emph{eṅkaḷ tokutikku oru putiya maruttuvamaṉai vēṇṭum.} — Our constituency needs a new hospital.
  \item \emph{Recall:} say the Tamil for a candidate, then the Tamil for to vote, then say \emph{tokuti} again
\end{itemize}

\begin{tcolorbox}[breakable,colback=teal!4,colframe=teal!35!black,title={Before you move on}]
What does \ta{தொகுதி} mean? (\textquotedblleft{}A constituency; a volume\textquotedblright{}.) Say it once more.
\end{tcolorbox}
\section[\texorpdfstring{\ta{பிரச்சாரம்} (piraccāram)}{பிரச்சாரம் (piraccāram)}]{\ta{பிரச்சாரம்} (piraccāram) — a campaign}

\label{lesson:TA-C374-piraccaram}



\begin{quote}
Before the new one: say the Tamil for to vote, then the Tamil for a constituency.
\end{quote}

\subsection*{You'll want to know: \ta{பிரச்சாரம்}}
\textbf{\ta{பிரச்சாரம்}} — \emph{piraccāram} — \textquotedblleft{}a campaign\textquotedblright{}.

\begin{quote}
\ta{வேட்பாளர்கள்} \ta{வீடு} \ta{வீடாகப்} \ta{பிரச்சாரம்} \ta{செய்தார்கள்}.
\end{quote}

\begin{quote}
\emph{vēṭpāḷarkaḷ vīṭu vīṭākap piraccāram ceytārkaḷ.}
\end{quote}

\begin{quote}
The candidates campaigned from house to house.
\end{quote}

\ta{பிரச்சாரம்} \ta{செய்}, to campaign.

\begin{cousinweb}
\textbf{\ta{பிரச்சாரம்}} (\emph{piraccāram}) is borrowed from Sanskrit \emph{pracāra}, \textquotedblleft{}coming forth, currency, prevalence\textquotedblright{} (Monier-Williams): spreading a message about.
\end{cousinweb}

\begin{grammarlens}[title={What you've built}]
One more word for giving reasons and opinions.
\end{grammarlens}

\subsection*{Your turn}
\begin{itemize}
\raggedright
  \item \emph{Say it:} \emph{piraccāram}
  \item \emph{Say it:} \emph{piraccāram}, once more
  \item \emph{Say it:} \emph{vēṭpāḷarkaḷ vīṭu vīṭākap piraccāram ceytārkaḷ.} — The candidates campaigned from house to house.
  \item \emph{Recall:} say the Tamil for to vote, then the Tamil for a constituency, then say \emph{piraccāram} again
\end{itemize}

\begin{tcolorbox}[breakable,colback=teal!4,colframe=teal!35!black,title={Before you move on}]
What does \ta{பிரச்சாரம்} mean? (\textquotedblleft{}A campaign\textquotedblright{}.) Say it once more.
\end{tcolorbox}
